\clearpage\fancyhead[L]{\small PALLAS.}\hypertarget{native-325-lower}{}\label{native-325-lower}
\begin{center}\textbf{PALLAS.}\par\vspace{4pt}\hrule\vspace{4pt}\fontsize{10}{12}\selectfont Monatliche Correspondenz. Vol. XXII. pp. 591\,..\,595. December 1810.\par\vspace{4pt}\hrule\vspace{4pt}\end{center}

\par\indent “I hasten to share some important news with you, my dearest friend,” wrote Professor \textsc{Gauss}. “I wrote to you recently that I was calculating the perturbations that Pallas had undergone owing to Jupiter in the years 1803\,..\,1811. In the paper I have mentioned to you several times, I showed how poorly purely elliptic elements can represent all six oppositions observed so far. I finished the work on those perturbations a few days ago. It was therefore immensely satisfying to find that, once those perturbations were taken into account, all six oppositions could be reconciled more closely than I had expected. Here is the table of discrepancies:

\begin{center}\fontsize{10}{12}\selectfont\setlength{\tabcolsep}{4pt}\renewcommand{\arraystretch}{1.15}\begin{tabular}{r|r|r}
 & \shortstack{Mean\\longitude} & \shortstack{Heliocentr.\\latitude}\\\hline
1803 & $+1''3$ & $-1''0$\\
1804 & $-3.8$ & $+4.1$\\
1805 & $+3.9$ & $+6.0$\\
1807 & $-3.3$ & $+3.9$\\
1808 & $+3.2$ & $-17.0$\\
1809 & $-1.4$ & $-3.9$\\
\end{tabular}\end{center}
\noindent (As is well known, the latitude for 1808 was determined very poorly.)

\SourceLeaf{326}\par\indent “Since the elements had changed \emph{very substantially} at the heliocentric conjunction with Jupiter in 1809, I was curious to see how far off my ephemeris in the October issue of the \emph{Monatl. Corresp.} might be. So I used my new results to calculate a position, taking the perturbations into account. For noon in Göttingen on 9 January 1811, I found

{\fontsize{10}{12}\selectfont\[\begin{array}{llr|l}\text{AE }\PallasSymbol & 151^\circ\;37' & \text{thus }1^\circ\;40'\text{ less} & \raisebox{-0.6\baselineskip}{than in the ephemeris.}\\\text{Decl.} & 22^\circ\;10'\text{ S} & \text{—}\quad 11'\text{ S} & \end{array}\]}

\noindent I have just calculated a second position, for noon in Göttingen on 26 October 1810

{\fontsize{10}{12}\selectfont\[\begin{array}{llr|l}\text{AE }\PallasSymbol & 134^\circ\;55' & \text{thus }45'\text{ less} & \raisebox{-0.6\baselineskip}{than in the ephemeris}\\\text{Decl.} & 14^\circ\;28'\text{ S} & \text{—}\quad 6'\text{ S} & \end{array}\]}

\Needspace{15\baselineskip}\noindent If you would like to calculate individual positions yourself, you can use the following elements:

\begin{center}\fontsize{10}{12}\selectfont\setlength{\tabcolsep}{4pt}\renewcommand{\arraystretch}{1.15}\begin{tabular}{llr}
Epoch 1811 Göttingen & .\;.\;.\;.\;.\;. & $126^\circ\;32'\;52''$\\
daily mean trop. motion & .\;. & $769''012$\\
Logarithm of the semimajor axis & .\;.\;.\;. & $0.44277$\\
Eccentricity & .\;.\;.\;.\;.\;.\;.\;.\;.\;. & $0.24162\ (=\operatorname{Sin}13^\circ\;18'\;56'')$\\
Perihelion for the above epoch & .\;.\;.\;.\;. & $120^\circ\;55'\;5''$\\
Node & .\;.\;.\;.\;.\;.\;.\;.\;.\;.\;.\;. & $172\;32\;30$\\
Inclination & .\;.\;.\;.\;.\;.\;.\;.\;.\;.\;. & $34\;35\;55$\\
\end{tabular}\end{center}

\noindent Before sending this, I can add one more position, for noon in Göttingen on 15 Dec. 1810.

{\fontsize{10}{12}\selectfont\[\begin{array}{llr|l}\text{AE }\PallasSymbol & 149^\circ\;37' & \text{thus }1^\circ\;15'\text{ less} & \raisebox{-0.6\baselineskip}{than in the ephemeris}\\\text{Decl.} & 21''\;37'\text{ S} & \text{—}\quad 12'\text{ farther south} & \end{array}\]}

\par\indent “I now ask you to use the first clear night to look for Pallas. It has been impossible here so far. You can easily imagine how impatiently I await confirmation of this remarkable result, which has never occurred before with the new planets.

\par\indent “The correction to the ephemeris in the October issue, obtained by interpolation, is as follows:
\Needspace{12\baselineskip}
\begin{center}\fontsize{10}{12}\selectfont\setlength{\tabcolsep}{4pt}\renewcommand{\arraystretch}{1.15}\begin{tabular}{ll|r|l}
 &  & AE. & Decl.\\\hline
1810, 26 October & Noon & $-45'$ & $6'$ S\\
--- 7 November & Midnight & 49 & 8\\
--- 20 November & Noon & 56 & 10\\
--- 2 December & Midnight & 65 & 11\\
--- 15 December & Noon & 75 & 12\\
--- 27 December & Midnight & 87 & 12\\
1811, 9 January & Noon & 101 & 11\\
\end{tabular}\end{center}

\par\indent Such remarkable results and the wish of our most esteemed friend were enough to prompt us to search for Pallas as diligently as we could. Unfortunately, the sky had been extremely unfavorable for six weeks, and this frustrated almost all our efforts. On 18 December it cleared a little. At two in the morning we saw three stars in the field of the telescope of the parallactic instrument. One of them was very probably Pallas, but a definite observation became impossible when the sky soon clouded over again. We could find no position for any of those three stars, either in \emph{Piazzi} or in the \emph{Histoire céleste}. On 19 De\SourceLeaf{327}cember we finally made a meridian observation. Exactly where the corrected ephemeris placed Pallas, we observed a star of magnitude 7...8 that was very probably Pallas.

\begin{center}\fontsize{10}{12}\selectfont\setlength{\tabcolsep}{4pt}\renewcommand{\arraystretch}{1.15}\begin{tabular}{c|c|c|c}
1809 & \shortstack{mean\\Seeberg time} & AR. app. $\PallasSymbol$ & \shortstack{Southern\\declination}\\\hline
19 December & $16^{\text{u}}\;9^{\text{m}}\;7^{\text{s}}3$ & $150^\circ\;20'\;13''5$ & $22^\circ\;58''7$\\
\end{tabular}\end{center}

\par\indent The declination could of course be observed only with the transit instrument and is uncertain at the level of a minute. The right ascension, however, is accurate and agrees exceedingly well with the corrected determination. Pallas was quite invisible in the telescope of the quadrant. A small star close to Pallas, whose position we found to be
\[\text{AR. }149''\;28'\;36''.\quad\text{Southern decl. }21^\circ\;59'\]
\noindent had remained in exactly the same position. We hesitated to tell Professor \textsc{Gauss} about this single observation, hoping that a second one would soon let us confirm that it really was Pallas. Unfortunately the sky has denied us that pleasure so far, so we are sending him our presumed position for Pallas. \textsc{Gauss}, who shares our conjecture, wrote back as follows:

\par\indent “I am very pleased that my calculation of the perturbations of Pallas has paid off so well. The weather has so far made it impossible for me to look for the planet, but I no longer doubt that the star you observed was Pallas. \textsc{Harding} has no star at that spot on his chart, but he has marked the other one that you also saw on 18 Dec., although it does not appear in the \emph{Histoire céleste}.”
